% Literature review (Phase 1). Citation keys resolve against reports/phase1/refs.bib.
% Papers gathered in the other section issues are merged in when the report is assembled.

\section{Literature Review}
\label{sec:literature_review}

\subsection{Prognostics and Remaining Useful Life}
Condition-based maintenance schedules maintenance from the measured condition of a
machine rather than from a fixed calendar, and prognostics, estimating how long a
component can keep operating, is its forward-looking half~\cite{Jardine2006review}.
The quantity to estimate is the remaining useful life (RUL): the number of operating
cycles left before failure. Reviews group RUL methods into physics-based models,
which need an explicit model of the failure mechanism; data-driven models, which learn
degradation from condition-monitoring data; and hybrids of the two~\cite{Lei2018machinery}.
Statistical data-driven approaches, from regression on degradation signals to
stochastic-process models, are surveyed by Si et al.~\cite{Si2011remaining}.
Lei et al.~\cite{Lei2018machinery} divide the data-driven pipeline into data
acquisition, health-indicator construction, health-stage division and RUL prediction.
The last two steps motivate our reporting of errors separately for the early, middle
and late stages of degradation.

\subsection{The C-MAPSS Benchmark}
Run-to-failure data from real engines is scarce, so simulation fills the gap.
Saxena et al.~\cite{Saxena2008damage} generated turbofan degradation data with the
C-MAPSS simulator. Each engine starts from a random, unknown level of initial wear; a
fault then reduces the flow and efficiency of engine modules at an exponential rate,
and the engine fails when a health index built from its operating margins reaches zero.
The resulting data set~\cite{Saxena2008turbofan} contains four subsets, FD001 to FD004,
that combine one or six operating conditions with one or two fault modes. For test
engines it gives only a truncated trajectory and the true RUL at its last cycle.
Saxena et al.~\cite{Saxena2008damage} also defined the asymmetric score used in the
PHM08 challenge, which penalises late predictions more heavily than early ones because
an unexpected failure costs more than early maintenance; a broader set of prognostic
metrics is discussed in~\cite{Saxena2008metrics}. Six years later, Ramasso and
Saxena~\cite{Ramasso2014performance} reviewed more than seventy publications on these
data. They found that the absence of benchmark results, together with misunderstandings
about how the subsets relate, made reported results hard to compare, and they named
sensor noise, changing operating conditions and multiple fault modes as the main
difficulties of the benchmark.
